\documentclass[a4paper,10pt]{article}
\usepackage[left=1cm,top=3cm,right=2.5cm,bottom=4cm]{geometry}
\usepackage[utf8]{inputenc}
\usepackage[activeacute,spanish]{babel}
\usepackage{multirow}	% Para poder unir filas en las tablas
\usepackage{colortbl}	% Para colorear tablas
\usepackage{latexsym,amssymb,oldlfont,amsmath}
\usepackage{graphicx,graphics}
\usepackage{multicol}
\usepackage{color}
\usepackage{listings}
\usepackage{lastpage}
\usepackage{enumitem}



%\usepackage{fontspec}
%\setmainfont{Times New Roman}
\newcommand{\R}{{\mathbb R}}
\newcommand{\si}{{\mathcal S}^n}
\renewcommand{\theequation}{\thesection.\arabic{equation}} % Cambia formato salida numeracion equation
\makeatletter\@addtoreset {equation}{section} \makeatother % Tambien se usa
\renewcommand{\baselinestretch}{1.3}
\usepackage{amssymb}
\usepackage{hyperref}
\definecolor{midnightblue}{rgb}{0.1,0.1,0.44}


\newcommand{\dis}{\displaystyle}
\renewcommand{\theenumii}{\arabic{enumii}}

\def\Z{\mathbb{Z}}
\def\N{\mathbb{N}}
\def\R{\mathbb{R}}
\def\M{\mathcal{M}_n(\mathbb{R})}
\def\sin{\operatorname{sen}}
\def\e{\operatorname{e}}
\DeclareMathOperator{\arcsec}{arcsec}
\DeclareMathOperator{\arccot}{arccot}
\DeclareMathOperator{\arccsc}{arccsc}



\usepackage{etoolbox,fancyhdr,xcolor}
\fancyhf{} % sets both header and footer to nothing
\renewcommand{\headrulewidth}{0pt}
%\renewcommand{\footrulewidth}{0.5pt}
\pagestyle{fancy}
\headheight=60pt %para cambiar el tamaño del encabezadox
\fancyhead[L] %la "L" indica a la izquierda
{
\begin{minipage}{6cm}
\includegraphics[width=0.9\textwidth]{logo_udp.eps}
\end{minipage}
}

%\fancyhead[C]
%{
%\textbf{Cálculo II - Guía 1}
%}


%\fancyhead[R] %la "L" indica a la izquierda
%{

%}
%\begin{minipage}{4cm}
%\small %tamaño de letra pequeño
%{
%\textsf
%{
%Universidad de Chile\\
%...}}
%\end{minipage}}



%\fancyfoot[C]{}
%{
%\begin{minipage}{10cm}
%\color{midnightblue}
%\begin{center}Casa Matriz: Monjitas 527 Piso 15 - Santiago de Chile\\
%Fono (02)24367255 \end{center}
%\end{minipage}
%}
\fancyfoot[C]{}
%Página \thepage\ de \pageref{LastPage}}% Páginas Pares numeradas a la derecha
%\fancyfoot[LO]{\thepage} % Páginas Impares numeradas a la izquierda

\title{C2_Ay1_Ej5}

\begin{document}

\begin{center}
Cálculo II\\
Ayudantía Nº1 - Ejercicio Nº5\\
Primer Semestre 2017
\end{center}





5. Escriba $\displaystyle \sin^{n}x\cos^{m}x = \sin^{n-1}x\cos^{m}x\sin x$, y recordando que $\displaystyle \cos^{2}x = 1-\sin^{2}x$ demuestre la f\'{o}rmula de reducci\'{o}n
$$\int \cos^{m}x\sin^nxdx=-\frac{\cos^{m+1}x\sin^{n-1}x}{m+n}+\frac{n-1}{m+n}\int \cos^mx\sin^{n-2}xdx$$
Luego, \'{u}sela para calcular $\dis \int \sin^{4}x\cos^2x dx$


\vspace{5mm}

Solución:

\begin{align*}
 I &= \int \sin^{n}(x)\cos^{m}(x)\, dx \\
 	I	&= \int \sin^{(n-1)}(x)\cos^{m}(x)\sin(x) \, dx
\end{align*}

Usando integración por partes:
\begin{align*}
 u  &=  (\sin (x))^{(n-1)}  &   v  &= -\frac{(\cos(x))^{(m+1)}}{m+1} \\
 du &= (n-1)(\sin (x))^{(n-2)} \cos(x)  \, dx     &   dv &= (\cos(x))^m\sin(x) dx
\end{align*}

Obtenemos lo siguiente:
\begin{align*}
 I	&= \displaystyle \sin^{(n-1)}(x)\cos^{m}(x)\sin(x) \, dx  \\
 	I	&= \displaystyle  -\frac{(\cos(x))^{(m+1)}}{m+1}(\sin (x))^{(n-1)} + \frac{n-1}{m+1} \int (\cos(x))^{(m+2)}(\sin (x))^{(n-2)}  \, dx \\
    I	&= \displaystyle  -\frac{(\cos(x))^{(m+1)}}{m+1}(\sin (x))^{(n-1)} + \frac{n-1}{m+1} \int (\cos(x))^m(\sin (x))^{(n-2)} {(\cos (x))}^2  \, dx
\end{align*}

Utilizando $\displaystyle \cos^{2}x = 1-\sin^{2}x$ obtenemos lo siguiente:
\begin{align*}
    \int \sin^{n}(x)\cos^{m}(x)\, dx	&= \displaystyle  -\frac{(\cos(x))^{(m+1)}}{m+1}(\sin (x))^{(n-1)} + \frac{n-1}{m+1} \int (\cos(x))^m(\sin (x))^{(n-2)} {(\cos (x))}^2  \, dx \\
    \int \sin^{n}(x)\cos^{m}(x)\, dx	&= \displaystyle  -\frac{(\cos(x))^{(m+1)}}{m+1}(\sin (x))^{(n-1)} + \frac{n-1}{m+1} \int (\cos(x))^m(\sin (x))^{(n-2)} (1-{(\sin (x))}^2)  \, dx\\
    \int \sin^{n}(x)\cos^{m}(x)\, dx	&= \displaystyle  -\frac{(\cos(x))^{(m+1)}}{m+1}(\sin (x))^{(n-1)} + \frac{n-1}{m+1} (\int (\cos(x))^m(\sin (x))^{(n-2)}\, dx -\int \sin^{n}(x)\cos^{m}(x)  \, dx)\\
    \int \sin^{n}(x)\cos^{m}(x)\, dx	&= \displaystyle  -\frac{(\cos(x))^{(m+1)}}{m+1}(\sin (x))^{(n-1)} + \frac{n-1}{m+1} \int (\cos(x))^m(\sin (x))^{(n-2)}\, dx -\frac{n-1}{m+1}\int \sin^{n}(x)\cos^{m}(x)  \, dx
\end{align*}

\begin{align*}
    \int \sin^{n}(x)\cos^{m}(x)\, dx +\frac{n-1}{m+1}\int \sin^{n}(x)\cos^{m}(x)  \, dx &= \displaystyle  -\frac{(\cos(x))^{(m+1)}}{m+1}(\sin (x))^{(n-1)} + \frac{n-1}{m+1} \int (\cos(x))^m(\sin (x))^{(n-2)}\, dx \\
    (1+ \frac{n-1}{m+1})\int \sin^{n}(x)\cos^{m}(x)\, dx &= \displaystyle  -\frac{(\cos(x))^{(m+1)}}{m+1}(\sin (x))^{(n-1)} + \frac{n-1}{m+1} \int (\cos(x))^m(\sin (x))^{(n-2)}\, dx\\
    (\frac{m+n}{m+1})\int \sin^{n}(x)\cos^{m}(x)\, dx &= \displaystyle  -\frac{(\cos(x))^{(m+1)}}{m+1}(\sin (x))^{(n-1)} + \frac{n-1}{m+1} \int (\cos(x))^m(\sin (x))^{(n-2)}\, dx
\end{align*}

\begin{align*}
    \int \sin^{n}(x)\cos^{m}(x)\, dx &= (\frac{m+1}{m+n})(\displaystyle  -\frac{(\cos(x))^{(m+1)}}{m+1}(\sin (x))^{(n-1)} + \frac{n-1}{m+1} \int (\cos(x))^m(\sin (x))^{(n-2)}\, dx)\\
     \int \sin^{n}(x)\cos^{m}(x)\, dx &= \displaystyle  -\frac{(\cos(x))^{(m+1)}}{m+n}(\sin (x))^{(n-1)} + \frac{n-1}{m+n} \int (\cos(x))^m(\sin (x))^{(n-2)}\, dx
\end{align*}
Quedando demostrada la formula de reducción.

\vspace{10mm}

Luego calculamos $\dis \int \sin^{4}x\cos^2x dx$

\begin{align*}
 \dis \int \sin^{4}x\cos^2x dx	&= \displaystyle -\frac{\cos^{3}(x)\sin^{3}(x)}{6} + \frac{3}{6} \int \cos^{3}(x)\sin^{3}(x) \, dx  \\
 	\dis \int \sin^{4}x\cos^2x dx	&= \displaystyle -\frac{\cos^{3}(x)\sin^{3}(x)}{6} + \frac{1}{2}( -\frac{\cos^{3}(x)\sin(x)}{4} + \frac{1}{4} \int \cos^{2}(x) \, dx ) \\
    \dis \int \sin^{4}x\cos^2x dx	&= \displaystyle -\frac{\cos^{3}(x)\sin^{3}(x)}{6} - \frac{1}{2}\frac{\cos^{3}(x)\sin(x)}{4} + \frac{1}{8} (\frac{x}{2} + \frac{1}{4}\sin(2x) +C) \\
   \dis \int \sin^{4}x\cos^2x dx	&= \displaystyle -\frac{\cos^{3}(x)\sin^{3}(x)}{6} - \frac{\cos^{3}(x)\sin(x)}{8} + \frac{x}{16} + \frac{1}{32}\sin(2x) +C
\end{align*}
\end{document}
